% TOP_PROBLEM: {"id": 3121, "title": "Minimal graphs with a prescribed number of spanning trees", "category_id": 3, "category": {"id": 3, "name": "graph_theory", "display_name": "Graph Theory", "description": "Problems involving graphs, networks, and their properties.", "slug": "graph-theory", "order_index": 3, "created_at": "2026-07-31T15:26:25.671Z"}, "status": "open", "problem_number": "OPG-37420", "set_id": 12}
% FIRST_POSTED: 2026-10-01
% ENTRY_KIND: statement_only
% UnsolvedMath ID 3121; source catalogue code OPG-37420
% !TeX program = lualatex
% Definition and sources only; this file is not an LLM solution attempt.
\documentclass[11pt,a4paper]{article}
\usepackage[margin=25mm]{geometry}
\usepackage{fontspec}
\setmainfont{lmroman10-regular.otf}[BoldFont=lmroman10-bold.otf,
 ItalicFont=lmroman10-italic.otf,BoldItalicFont=lmroman10-bolditalic.otf]
\usepackage{amsmath,amssymb,mathtools}
\usepackage{unicode-math}
\setmathfont{latinmodern-math.otf}
\AtBeginDocument{\let\setminus\smallsetminus}
\usepackage{xurl,hyperref}
\hypersetup{colorlinks=true,urlcolor=blue}
\setcounter{secnumdepth}{0}
\setlength{\parindent}{0pt}
\setlength{\parskip}{5pt}
\setlength{\emergencystretch}{3em}
\begin{document}
\section*{Minimal graphs with a prescribed number of spanning trees}\label{3121}
{\small UnsolvedMath ID 3121 · OPG-37420 · Graph Theory · OpenGarden}
\subsection{Definitions and mathematical statement}
For a finite simple undirected graph $G$, a spanning tree is a connected
acyclic subgraph containing every vertex. Let $\tau(G)$ be the number
of its spanning trees, with distinct edge sets counted separately.
For each integer $n\ge3$, define
\[
 \alpha(n)=\min\{|V(G)|:G\text{ finite, simple and undirected},\
                                  \tau(G)=n\}.
\]
The minimum exists: deleting any one edge of the cycle $C_n$ gives
one of its exactly $n$ spanning trees, so $\alpha(n)\le n$.
Any graph contributing to the minimum is connected because its tree
count is positive. Loops, parallel edges and directed arcs are not allowed.

The conjecture asserts
\[
 \lim_{n\to\infty}\frac{\alpha(n)}{\log n}=0,
\]
where $n$ tends to infinity through all integers and $\log$ may be taken
to be the natural logarithm. Equivalently, for every $\varepsilon>0$
there exists $N_\varepsilon$ such that every integer $n\ge N_\varepsilon$
is the exact spanning-tree count of a simple graph on at most
$\varepsilon\log n$ vertices.

The number $n$ here is the prescribed count, not the number of graph
vertices. The conclusion is uniform over all sufficiently large counts;
constructions only for a sparse subsequence of highly composite integers
do not suffice. Nor is it enough to find a graph with at least $n$
spanning trees. The exact equality is what makes the minimization an
arithmetic realization problem.
\subsection{Short English statement}
Can every sufficiently large integer be realized as an exact spanning-tree count using a number of vertices negligible compared with its logarithm?
\subsection{Sources}
\begin{itemize}
\item[S1] ULAM.AI and the UnsolvedMath Contributors, supplied problems.json, record 3121 (OPG-37420), OpenGarden collection. Catalogue identity and source transcription; CC BY 4.0.
\url{https://huggingface.co/datasets/ulamai/UnsolvedMath}.
\item[S2] Open Problem Garden, Minimal graphs with a prescribed number of spanning trees. Original problem and discussion; the mathematical formulation below is expanded from the supplied source record.
\url{http://www.openproblemgarden.org/op/minimal_graphs_with_a_prescribed_number_of_spanning_trees}.
\end{itemize}
\end{document}
